\section{Introduction}\label{sec:introduction}

Building models that solve a diverse set of tasks has become a dominant paradigm in the domains of vision and language. In natural language processing, large pretrained models have demonstrated remarkable zero-shot learning of new language tasks~\citep{brown2020language,chowdhery2022palm,hoffmann2022training}. Similarly, in computer vision, models such as those proposed in~\citep{radford2021learning,alayrac2022flamingo}
have shown remarkable zero-shot classification and object recognition capabilities. A natural next step is to use such tools to construct agents that can complete different decision making tasks across many environments.

However, 
training such agents faces the inherent challenge of environmental diversity, since different environments operate with distinct state action spaces (e.g., the joint space and continuous controls in MuJoCo are fundamentally different from the image space and discrete actions in Atari). Such diversity hampers knowledge sharing, learning, and generalization across tasks and environments. Although substantial effort has been devoted to encoding different environments with universal tokens in a sequence modeling framework~\citep{reed2022generalist}, it is unclear whether such an approach can preserve the rich knowledge embedded in pretrained vision and language models and leverage this knowledge to transfer to downstream reinforcement learning~(RL) tasks. Furthermore, it is difficult to construct reward functions that specify different tasks across environments.

In this work, we address the challenges in environment diversity and reward specification by leveraging \emph{video} (i.e., image sequences) as a universal interface for conveying action and observation behavior in different environments, and \emph{text} as a universal interface for expressing task descriptions. In particular, we design a video generator as a planner that sequentially conditions on a current image frame and a text passage describing a current goal (i.e., the next high-level step) to generate a trajectory in the form of an image sequence, after which an inverse dynamics model is used to extract the underlying actions from the generated video. 
Such an approach allows the universal nature of language and video to be leveraged in generalizing to novel goals and tasks across diverse environments. Specifically, we instantiate the text-conditioned video generation model using video diffusion.
A set of underlying actions are then regressed from the synthesized frames and used to construct a policy to implement the planned trajectory. Since the language description of a task is often highly correlated with control actions, text-conditioned video generation naturally focuses generation on action relevant parts of the video. The proposed model, \textit{\model}, is visualized in \fig{fig:overview}. 

We have found that formulating policy generation via text-conditioned video synthesis yields the following advantages:
    
\myparagraph{Combinatorial Generalization.}  The rich combinatorial nature of language can be leveraged to synthesize novel combinatorial behaviors in the environment. This enables the proposed approach to rearrange objects to new unseen combinations of geometric relations, as shown in \sect{sect:combinatorical}.

\myparagraph{Multi-task Learning.} Formulating action prediction as a video prediction problem readily enables learning across many different tasks. We illustrate in \sect{sect:multi} how this enables learning across language-conditioned tasks and generalizing to new ones at test time  without finetuning.

\input{figText/overview.tex}
\myparagraph{Action Planning.} The video generation procedure corresponds to a planning procedure where a sequence of frames representing actions is generated to reach the target goal. 
Such a planning procedure is naturally hierarchical: a temporally sparse sequence of images toward a goal can first be generated, before being refined with a more specific plan. %
Moreover, the planning procedure is %
steerable, in the sense that the plan generation can be biased by new constraints introduced
at test-time through test-time sampling. 
Finally, plans are produced in a video space that is naturally interpretable by humans, making action verification and plan diagnosis easy. 
We illustrate the efficacy of hierarchical sampling in \tbl{tbl:ablations} and steerability in \fig{fig:comb_adapt}.

\myparagraph{Internet-Scale Knowledge Transfer.} By pretraining a video generation model on a large-scale text-video dataset recovered from the internet, one can recover a vast repository of ``demonstrations'' that aid the construction of a text-conditioned policy in novel environments. We illustrate how this enables the realistic synthesis of robot motion videos from language instructions in \sect{sect:real}.

The main contribution of this work is to formulate text-conditioned video generation as a universal planning strategy from which diverse behaviors can be synthesized. While such an approach departs from typical policy generation in RL, where subsequent actions to execute are directly predicted from a current state, we illustrate that \model exhibits notable generalization advantages over traditional policy generation %
methods
across a variety of  domains.

